\section{A compact follow-up research program}

The calculations suggest three tightly linked projects:
\begin{enumerate}
\item \textbf{Formalize the $d\ge3$ family.} Prove the computational/Fourier result in full generality, then verify or correct the arbitrary-MUB adaptation.
\item \textbf{Resolve the qubit case.} Use the six-moment reduction to formulate a constrained minimization or semidefinite-relaxation proof. A counterexample search should focus outside every known proved class.
\item \textbf{Sharpen replacement bounds.} Determine the best bound as a function of dimension, number and geometry of measurement bases, marginal entropies, global entropy or purity, and perhaps local support dimensions.
\end{enumerate}

\section{Current answers to the motivating questions}

\begin{tcolorbox}[colback=cqcpale,colframe=cqcblue!60!black,boxrule=0.55pt,arc=1.5pt,left=6pt,right=6pt,top=5pt,bottom=5pt]
\textbf{Higher dimensions.} Yes: the computational/Fourier family gives an explicit counterexample in every $d\times d$ system with $d\ge3$, with closed-form gap $\Delta_d(\eps)=\hbin(\eps)-2\hbin(\eps/d)$. The same compression argument appears to adapt the construction to any prescribed local MUB pair.

\textbf{Two qubits.} No two-qubit counterexample is known from this work. This higher-dimensional family has an exact no-go at $d=2$, and exploratory searches have found no alternative. The unrestricted qubit case remains open and may survive as a separate theorem.

\textbf{Replacement bounds.} Conditional-entropy, mixedness, purity, outcome-distribution, and basis-overlap bounds remain valid and provide a structured replacement for the failed universal inequality.
\end{tcolorbox}

\section*{Reproducibility}

The qutrit exact certificate, all-dimension verifier, and exploratory two-qubit search are in the public repository:
\begin{itemize}
\item \path{scripts/verify_cqc_exact.py};
\item \path{scripts/verify_cqc_dimension_family.py};
\item \path{scripts/search_cqc_qubit.py}.
\end{itemize}
The exact qutrit result and the algebraic formulas above do not depend on the numerical searches.

\begin{thebibliography}{9}
\bibitem{Schneeloch2014}
J. Schneeloch, C. J. Broadbent, and J. C. Howell,
``Uncertainty relation for mutual information,''
\emph{Physical Review A} \textbf{90}, 062119 (2014).
\href{https://doi.org/10.1103/PhysRevA.90.062119}{doi:10.1103/PhysRevA.90.062119}.

\bibitem{Iqbal2026}
H. Iqbal,
``On the CQC conjecture: a sufficient condition and an extension,''
\emph{Quantum Information Processing} \textbf{25}, 250 (2026).
\href{https://doi.org/10.1007/s11128-026-05258-2}{doi:10.1007/s11128-026-05258-2}.

\bibitem{Berta2010}
M. Berta, M. Christandl, R. Colbeck, J. M. Renes, and R. Renner,
``The uncertainty principle in the presence of quantum memory,''
\emph{Nature Physics} \textbf{6}, 659--662 (2010).
\href{https://doi.org/10.1038/nphys1734}{doi:10.1038/nphys1734}.

\bibitem{ColesPiani2014}
P. J. Coles and M. Piani,
``Improved entropic uncertainty relations and information exclusion relations,''
\emph{Physical Review A} \textbf{89}, 022112 (2014).
\href{https://doi.org/10.1103/PhysRevA.89.022112}{doi:10.1103/PhysRevA.89.022112}.
\end{thebibliography}

\vfill
{\footnotesize\color{cqcgray}
This working note is part of the public record at
\url{https://github.com/EthanOConnor/cqc-qutrit-counterexample}.
The repository's \texttt{PROVENANCE.md} gives the complete attribution and review-status statement.
}
